% -------------------------------------------------------------------------------------
\section{规范对称性与超软退耦}\label{sec:scet-decoupling}

在第~9~章的图形方法中，软胶子与喷注的耦合经程函近似与 Ward 恒等式归并为软 Wilson 线；在 SCET 中，同一结果来自一个场的重定义，它把超软胶子从领头阶共线拉氏量中整体消去。要做到这一点，共线拉氏量的每一部分——包括规范固定项与鬼场项——都必须对超软场协变，所以先把共线胶子部分写成以超软场为背景的形式。共线胶子与共线鬼场的领头阶拉氏量是 BPS 的 (24) 式\cite{BauerPirjolStewart2002}，写成坐标空间形式为
\begin{equation}\label{eq:scet-Lcg}
\begin{split}
\mathcal L^{(0)}_{cg}={}&\underbrace{\frac{1}{2g^2}\Tr\Big\{\big[\ii\mathcal D^\mu+gA_n^\mu,\ii\mathcal D^\nu+gA_n^\nu\big]\Big\}^2}_{\mathcal L_F}
+\underbrace{\frac1\alpha\Tr\Big\{\big[\ii\mathcal D_\mu,A_n^\mu\big]\Big\}^2}_{\mathcal L_{\rm gf}}\\
&+\underbrace{2\Tr\Big\{\bar c_n\big[\ii\mathcal D_\mu,\big[\ii\mathcal D^\mu+gA_n^\mu,c_n\big]\big]\Big\}}_{\mathcal L_{\rm gh}},
\end{split}
\end{equation}
其中 $\ii\mathcal D^\mu\equiv\frac{n^\mu}{2}\,\ii\bar n\cdot\partial+\ii\partial_\perp^\mu+\frac{\bar n^\mu}{2}\,\ii n\cdot D_{us}$，$\ii n\cdot D_{us}\equiv\ii n\cdot\partial+g\,n\cdot A_{us}(x_-)$，$\Tr\,t^at^b=\frac12\delta^{ab}$，$c_n$、$\bar c_n$ 为伴随表示的共线鬼场，$\alpha$ 为规范参数。由 $[\ii D_\mu,\ii D_\nu]=\ii gG_{\mu\nu}$，$\mathcal L_F=-\frac14G^a_{\mu\nu}G^{a\mu\nu}$，$\mathcal L_{\rm gf}=-\frac{1}{2\alpha}(\mathcal D\cdot A_n)^a(\mathcal D\cdot A_n)^a$，$\mathcal L_{\rm gh}=-\bar c_n^a(\mathcal D^\mu D_{\mu}c_n)^a$，其中 $\mathcal D$ 是只含背景场 $A_{us}$ 的伴随协变导数，$D=\mathcal D-\ii g[A_n,\cdot\,]$；$g\to0$、$A_{us}\to0$ 时它们化为通常的自由胶子、规范固定与鬼场项。因此\cref{eq:scet-Lcg}是以 $A_{us}$ 为背景场、$A_n$ 为量子场的背景场规范 QCD 作用量的领头幂次\cite{BauerPirjolStewart2002}。% w18_s3 (g)
\emph{共线}规范变换 $U_c(x)$（导数与共线动量同阶，在无穷远处趋于 $1$）：$\xi_n\to U_c\xi_n$，$A_n\to U_cA_nU_c^\dagger+\frac1gU_c[\ii\mathcal D,U_c^\dagger]$，超软场不变；\emph{超软}规范变换 $V(x)$（导数 $\sim\lambda^2Q$，作用于共线场时取 $V(x_-)$）：$\xi_n\to V\xi_n$，$A_n\to VA_nV^\dagger$，$A_{us}\to VA_{us}V^\dagger+\frac1gV[\ii\partial,V^\dagger]$。二者都使 $\ii\mathcal D+gA_n\to U(\ii\mathcal D+gA_n)U^\dagger$（$U=U_c$ 或 $V$），超软变换还单独使 $\ii\mathcal D\to V\,\ii\mathcal D\,V^\dagger$。

\begin{proposition}[共线与超软规范对称性 \tagR]\label{prop:scet-gauge}
（$R_{\mathrm{def}}$：由算符定义直接推出。）把\cref{eq:scet-L0xi}与\cref{eq:scet-Lcg}看作场的表达式：
(a) $\mathcal L^{(0)}_{\xi\xi}$ 与场强项 $\mathcal L_F$ 在共线变换与超软变换下都不变；
(b) 规范固定项 $\mathcal L_{\rm gf}$ 与鬼场项 $\mathcal L_{\rm gh}$ 各自在超软变换（$c_n,\bar c_n\to Vc_nV^\dagger,V\bar c_nV^\dagger$）下不变，即对超软背景协变；二者都不是共线规范不变的（无论鬼场按伴随表示变换还是不变）；
(c) 令 BRST 变换作用于量子场，$sA_n=-\ii[\ii\mathcal D+gA_n,c_n]$（即 $sA^a_n=(Dc_n)^a$）、$sc_n^a=-\frac g2f^{abc}c_n^bc_n^c$、$s\bar c_n^a=B^a$、$sB^a=0$，超软场是不变的背景（\cref{def:prelim-brst}），则 $\mathcal L_{\rm gf}+\mathcal L_{\rm gh}$ 是 BRST 正合的，从而它与整个 $\mathcal L^{(0)}_{cg}$ 都是 BRST 不变的；单独的 $\mathcal L_{\rm gf}$ 或 $\mathcal L_{\rm gh}$ 不是。
（这订正了文献\cite[式~(24) 之后]{BauerPirjolStewart2002}关于鬼场项具有共线规范不变性的表述，见附录~Q 第~3~条。）（本命题只涉及有效拉氏量，不涉及强子态。）
\uses{def:scet-modes, prop:qcd-covariance, lem:prelim-brst-nilpotent, def:prelim-brst}
\end{proposition}

\begin{proof}
(a) 两类变换都使 $\ii\mathcal D^\mu+gA_n^\mu$ 按共轭变换（\cref{prop:qcd-covariance}，$\ii\partial$ 换成背景协变的 $\ii\mathcal D$），对易子随之共轭，迹不变，故 $\mathcal L_F$ 不变。$\mathcal L^{(0)}_{\xi\xi}$ 中的 $\ii n\cdot D$、$\ii\bar n\cdot D_c$、$\ii D_{\perp c}$ 正是 $\ii\mathcal D+gA_n$ 的三个分量，$\frac1{\ii\bar n\cdot D_c}\to U\frac1{\ii\bar n\cdot D_c}U^\dagger$，而 $\xi_n\to U\xi_n$、$U$ 与 Dirac 矩阵对易，故不变；对超软变换还用到 $V(x_-)$ 只依赖 $\bar n\cdot x$，$\bar n\cdot\partial V=\partial_\perp V=0$。
(b) 超软变换下 $[\ii\mathcal D_\mu,A_n^\mu]\to V[\ii\mathcal D_\mu,A_n^\mu]V^\dagger$，鬼场项的双重对易子同样按共轭变换。共线变换下 $[\ii\mathcal D_\mu,A_n^\mu]$ 获得非齐次项 $\frac1g[\ii\mathcal D_\mu,U_c[\ii\mathcal D^\mu,U_c^\dagger]]\neq0$；鬼场项的外层对易子 $[\ii\mathcal D_\mu,\cdot\,]$ 不是共线协变的，其变分既不因 $c_n\to U_cc_nU_c^\dagger$ 也不因鬼场不变而消失（对随机 su(3) 场，相对变分为 $\order{1}$）。% w18_s3 (b),(c); App Q item 3
(c) 令 $G^a\equiv(\mathcal D\cdot A_n)^a$，取规范费米子 $\Psi=\bar c_n^a(G^a+\frac\alpha2B^a)$。背景不变而 $G$ 对 $A_n$ 线性，$sG^a=(\mathcal D^\mu D_\mu c_n)^a$，故（$s$ 从左作用）$s\Psi=B^aG^a+\frac\alpha2B^aB^a-\bar c_n^a(\mathcal D^\mu D_\mu c_n)^a$；消去 $B^a=-G^a/\alpha$，前两项为 $-\frac1{2\alpha}G^aG^a=\mathcal L_{\rm gf}$，末项即 $\mathcal L_{\rm gh}$。$s^2=0$ 的证明与 QCD 相同（\cref{lem:prelim-brst-nilpotent}；背景不参与 $s$），故 $s(\mathcal L_{\rm gf}+\mathcal L_{\rm gh})=s^2\Psi=0$；$s$ 作用于 $(A_n,\xi_n)$ 是参数 $\propto c_n$ 的无穷小共线规范变换，由 (a)，$s\mathcal L_F=0$。单独看，$s\mathcal L_{\rm gh}=-B^a(\mathcal D^\mu D_\mu c_n)^a$（用 $s(D_\mu c_n)=s^2A_{n\mu}=0$），而 $s(B^aG^a+\frac\alpha2B^aB^a)=+B^a(\mathcal D^\mu D_\mu c_n)^a$：两项只在求和中相消。
\end{proof}

由此，$\chi_n=W_n^\dagger\xi_n$ 与 $\mathcal B^\mu_{n\perp}\equiv\frac1g[W_n^\dagger\,\ii D^\mu_{\perp c}W_n]$ 是共线规范不变的组元，在超软变换下按 $V(\cdots)$ 与 $V(\cdots)V^\dagger$ 变换。定义超软 Wilson 线（多极展开后的场）
\begin{equation}\label{eq:scet-Yn}
\begin{gathered}
Y_n(x)\equiv W[x,x-\infty n]=\mathrm P\exp\Big(\ii g\int_{-\infty}^0\dd s\;n\cdot A_{us}\big((x+sn)_-\big)\Big),\\
Y_n^{(+)}(x)\equiv W[x,x+\infty n]\quad\text{（指向未来的线）};
\end{gathered}
\end{equation}
因为 $(x+sn)_-=(\bar n\cdot x+2s)\frac n2$，$Y_n$ 只依赖于 $\bar n\cdot x$，故 $\bar n\cdot\partial Y_n=\partial_\perp Y_n=0$。

\begin{theorem}[超软退耦（BPS 场重定义） \tagR]\label{thm:scet-decoupling}
（$R_{\mathrm{def}}$：由算符定义直接推出。）设 $\mathcal L^{(0)}_{\xi\xi}$、$\mathcal L^{(0)}_{cg}$ 如\cref{eq:scet-L0xi,eq:scet-Lcg}。
(a)（关键恒等式）作为算符，
\begin{equation}\label{eq:scet-key}
\ii n\cdot D_{us}\,Y_n=Y_n\,\ii n\cdot\partial,\qquad\text{即}\quad Y_n^\dagger\,\ii n\cdot D_{us}\,Y_n=\ii n\cdot\partial ;
\end{equation}
% VF-SCET-6
对 $Y_n^{(+)}$ 同样成立。
(b) 作场重定义
\begin{equation}\label{eq:scet-BPS}
\xi_n=Y_n\,\xi_n^{(0)},\qquad A_n^\mu=Y_n\,A_n^{(0)\mu}\,Y_n^\dagger,\qquad c_n=Y_n\,c_n^{(0)}\,Y_n^\dagger\quad(\bar c_n\ \text{同样}),
\end{equation}
则 $\mathcal L^{(0)}_{\xi\xi}[\xi_n,A_n,n\cdot A_{us}]+\mathcal L^{(0)}_{cg}[A_n,c_n,n\cdot A_{us}]=\mathcal L^{(0)}_{\xi\xi}[\xi^{(0)}_n,A^{(0)}_n,0]+\mathcal L^{(0)}_{cg}[A^{(0)}_n,c^{(0)}_n,0]$，且变换的 Jacobi 行列式为 $1$：在领头幂次新共线场与超软胶子完全不相互作用。
成立的条件是：(1) 共线鬼场同时被重定义；(2) 共线规范固定项对超软背景协变（用 $\ii\mathcal D$ 而非 $\ii\partial$）；(3) Jacobi 行列式为 $1$；(4) 只在领头幂次；(5) $Y_n$ 由多极展开后的 $n\cdot A_{us}(x_-)$ 构成。（本定理是领头阶 SCET 拉氏量的恒等式，不涉及强子态。）
\uses{prop:scet-gauge, prop:prelim-wilson-covariance}
\end{theorem}

\begin{proof}
(a) 由\cref{prop:prelim-wilson-covariance}，对上端点沿 $n$ 求导得 $n\cdot\partial\,Y_n=\ii g\,n\cdot A_{us}(x_-)Y_n$。于是对任意 $f$，
$\ii n\cdot D_{us}(Y_nf)=\ii(n\cdot\partial Y_n)f+Y_n\,\ii n\cdot\partial f+g\,n\cdot A_{us}Y_nf=Y_n\,\ii n\cdot\partial f$。$Y_n^{(+)}$ 的上端点也是 $x$，论证不变（对随机 su(3) 背景场数值求解路径有序指数，误差 $\sim10^{-9}$）。% w18_s2 KEY IDENTITY
(b) $V\equiv Y_n^\dagger$ 只依赖于 $\bar n\cdot x$，是超软规范变换。由\cref{prop:scet-gauge}(a)(b)，$\mathcal L^{(0)}_{\xi\xi}+\mathcal L^{(0)}_{cg}$ 作为 $(\xi_n,A_n,c_n,\bar c_n;n\cdot A_{us})$ 的函数在这种变换下不变；把它作用于宗量 $(Y_n\xi^{(0)}_n,Y_nA^{(0)}_nY_n^\dagger,Y_nc^{(0)}_nY_n^\dagger;n\cdot A_{us})$，共线场变为 $(\xi^{(0)}_n,A^{(0)}_n,c^{(0)}_n)$，超软场变为
$Y_n^\dagger\,n\cdot A_{us}Y_n+\frac1gY_n^\dagger[\ii n\cdot\partial,Y_n]=Y_n^\dagger\,n\cdot A_{us}Y_n-Y_n^\dagger\,n\cdot A_{us}Y_n=0$（用 (a)）。所以重定义\cref{eq:scet-BPS}是只作用于共线扇区的“$n\cdot A_{us}=0$ 规范”变换，超软场本身不变。条件 (1)(2) 是\cref{prop:scet-gauge}(b) 的超软不变性所需：用 $\ii\partial$ 的规范固定项或不重定义的鬼场会留下与 $n\cdot A_{us}$ 的耦合；条件 (5) 保证 $Y_n$ 是超软变换并与 $\ii\bar n\cdot\partial$、$\ii\partial_\perp$ 对易。$Y_n$ 幺正、幺模且不依赖于共线场：对费米子是逐点 $\det Y_n=1$ 的转动，对胶子与鬼场是伴随表示中的正交转动，故 Jacobi 行列式为 $1$（条件 (3)）。次领头的超软耦合（含 $A_{us\perp}$、$\bar n\cdot A_{us}$ 或超软夸克）不是 $n\cdot A_{us}$ 的形式，不能只靠 $Y$ 因子吸收\cite{BauerPirjolStewart2002}（条件 (4)）。
\end{proof}

\paragraph{地位。}作为领头阶 SCET 拉氏量的恒等式，\cref{thm:scet-decoupling}是 [R]（$R_{\mathrm{def}}$），对 $\as$ 的所有阶成立；作为关于 QCD 振幅的陈述，它还需要模式完备、零箱减除正确、Glauber 模式不起作用——这些是 [A] 层次的输入（\cref{rem:scet-assumptions}）。

在重定义之下 $\chi_n=W_n^\dagger\xi_n=Y_n\chi^{(0)}_n$，两喷注流化为
\begin{equation}\label{eq:scet-JY}
\bar\chi_n\,\Gamma\,\chi_{\bar n}=\bar\chi^{(0)}_n\,Y_n^\dagger\,\Gamma\,Y_{\bar n}\,\chi^{(0)}_{\bar n},
\end{equation}
其在 $|X_nX_{\bar n}X_{us}\rangle$ 与真空之间的矩阵元因子化为 $\langle X_n|\bar\chi^{(0)}_n|0\rangle\,\Gamma\,\langle X_{\bar n}|\chi^{(0)}_{\bar n}|0\rangle\,\langle X_{us}|Y_n^\dagger Y_{\bar n}|0\rangle$（色指标依次缩并）：这是推力因子化中 $H\,J\otimes J\otimes S$ 结构的来源（\cref{fig:scet-wilson}(b)）。

\begin{example}[BPS 场重定义的树图演示]\label{ex:scet-decoupling-tree}
取一个被吸收的超软胶子 $n\cdot A_{us}=n\cdot\epsilon\,t^a\ee^{-\ii k\cdot x_-}$，$n\cdot k\neq0$。由\cref{eq:scet-Yn}，$Y_n=1-g\,t^a\frac{n\cdot\epsilon}{n\cdot k+\ii0}\ee^{-\ii k\cdot x_-}+\order{g^2}$，$\ii n\cdot\partial Y_n=-g\,t^an\cdot\epsilon\,\ee^{-\ii k\cdot x_-}+\order{g^2}$ 恰与 $g\,n\cdot A_{us}$ 相消。振幅层次：由流产生的出射 $n$ 共线夸克吸收该胶子，按\cref{prop:scet-lagrangian}的 Feynman 规则（$(p-k)^2\simeq-\bar n\cdot p\,n\cdot k$，$\bar u_nP_{\bar n}=\bar u_n$）
\[
\bar u_n(p)\,\big(\ii g\,t^a\,n\cdot\epsilon\,\tfrac{\bar{\sla n}}{2}\big)\,\ii\frac{\sla n}{2}\frac{\bar n\cdot p}{(p-k)^2+\ii0}\,\Gamma=\bar u_n(p)\,\frac{g\,t^a\,n\cdot\epsilon}{n\cdot k-\ii0}\,\Gamma .
\]
重定义之后同一耦合来自流中的 $Y_n^\dagger$，它给出 $g\,t^a n\cdot\epsilon/(n\cdot k+\ii0)$，而指向未来的 $Y_n^{(+)\dagger}=W[x+\infty n,x]$ 给出 $g\,t^an\cdot\epsilon/(n\cdot k-\ii0)$，与 Feynman 规则逐项相同；两者之差 $-2\pi\ii\,\delta(n\cdot k)$ 只在 $n\cdot k=0$ 处有支撑。% w18_s5 (F)
\uses{thm:scet-decoupling, prop:scet-lagrangian}
\end{example}

\begin{remark}[$Y_n$ 的方向与 $\ii0$]\label{rem:scet-y-direction}
拉氏量层次的证明只用到 $n\cdot D_{us}Y=0$，对 $Y_n$ 与 $Y_n^{(+)}$ 同样成立；振幅中的 $\ii0$ 则取决于方向。\cref{ex:scet-decoupling-tree}表明，领头阶 SCET 的 Feynman 规则本身对出射的共线夸克给出指向未来的 $Y$ 线：方向由匹配外态确定\cite{Chay:2004zn}。Rothstein--Stewart 指出，S 矩阵元中超软线的方向由流与入射、出射态的插值场共同决定，对 $n$–$\bar n$ 产生两条线都伸向 $+\infty$\cite[第~10.1~节]{RothsteinStewart2016}；他们还指出，SCET$_{\rm I}$ 中超软 Wilson 线的方向由物理方向决定，因为这些图没有减除，而软与共线 Wilson 线的方向依赖被（Glauber）零箱减除抵消\cite[第~4.2~节]{RothsteinStewart2016}。对单圈推力这一差别无关紧要（虚图无标度，实辐射有 $n\cdot k,\bar n\cdot k>0$）；$e^+e^-$、DIS 与 Drell--Yan 的半球软函数直到 $\order{\as^2}$ 都相等，尽管个别振幅的虚部符号不同\cite{Kang:2015moa}。此处 $\lambda^2Q$ 的模式是超软的；超软退耦\emph{不}消去 Glauber 相互作用：重定义之后 Glauber 拉氏量中只出现伴随表示的超软线，$n$ 与 $\bar n$ 扇区仍然耦合\cite[第~5.1~节]{RothsteinStewart2016}（\cref{sec:scet-glauber}）。
\end{remark}
